\documentclass[12pt,a4paper]{ctexart}
\usepackage[a4paper,margin=2.4cm]{geometry}
\usepackage{amsmath,amssymb,amsthm,bm,mathtools}
\usepackage{physics}
\usepackage{hyperref}
\usepackage{booktabs}
\usepackage{enumitem}
\usepackage{xcolor}
\usepackage{array}
\usepackage{longtable}
\usepackage{graphicx}
\usepackage{float}
\usepackage{siunitx}
\hypersetup{colorlinks=true,linkcolor=blue,citecolor=blue,urlcolor=blue}
\setlist[enumerate]{leftmargin=2.2em}
\setlist[itemize]{leftmargin=2em}
\allowdisplaybreaks[4]

\title{基于 SDF 的压力场接触理论：\\
一份面向实现的正式推导版文档}
\author{}
\date{}

\newtheorem{definition}{定义}
\newtheorem{theorem}{定理}
\newtheorem{proposition}{命题}
\newtheorem{corollary}{推论}
\newtheorem{remark}{备注}

\begin{document}
\maketitle

\begin{abstract}
本文将你前面几份补充材料中的思路，连同 Pressure Field Contact（PFC）论文中的核心理论，整理为一套自洽的正式推导版本。本文的目标不是简单复述已有公式，而是回答四个更根本的问题：
\begin{enumerate}
    \item 若几何由有符号距离函数（signed distance function, SDF）给出，是否能够实现与 PFC 同类的压力场接触理论；
    \item 为什么“等压面 + traction 面积分”可以改写成“窄带体积分”，以及这一改写何时是正确的；
    \item 为什么 band mechanics 与 sheet representation 必须分层，不能再把厚带直接当成零厚度接触片；
    \item 若数值结果与理论不符，应优先从哪些一致性条件与离散误差来源着手排查。
\end{enumerate}

此外，本文补入了一节面向一般三维曲面的正式推导：在局部坐标系 $(u,v,s)$ 中，把“横向 patch 几何 + 法向一维 band 积分”的思想推广为 local-normal 3D accumulator，从而把前文平面对平面基准下的 per-column 结构上升为一套适用于一般三维接触片的统一表述。

全文采用一个清晰的立场：\textbf{用 SDF 来实现压力场接触是可行的；真正需要重建的不是“是否改用 SDF”这一方向，而是“如何把几何、压力、平衡面、体积分测度、零厚度接触片表示、局部法向积分器以及 benchmark 校准组织成一条闭合的理论链”。}
\end{abstract}

\tableofcontents

\section{引言：我们到底要计算什么}

设两个名义上刚性的物体发生接触。真正要计算的对象并不是“有没有碰撞”这样一个布尔判断，而是下面三件事：
\begin{enumerate}
    \item 接触区域上每一点有多大的法向压力；
    \item 这些局部压力合起来总共产生多大的合力；
    \item 这些局部压力相对某个参考点总共产生多大的合力矩。
\end{enumerate}

如果只把接触看成一个点，那么许多与接触面积有关的现象就无法表达，例如：
\begin{itemize}
    \item 偏心压缩时的接触力矩；
    \item 面接触、条带接触中的支撑区域变化；
    \item 抓取、滚动、擦转等依赖接触片分布的现象。
\end{itemize}

PFC 的核心思想正是：\textbf{不要直接把接触简化成一个点，而是先定义一张由两侧压力场共同决定的平衡接触面，然后在这张面上积分 traction。}

本文要做的，是把这一思想完整地移植到 SDF 几何表示下，并进一步说明如何从“零厚度等压面理论”走到“适合 OpenVDB/NanoVDB/Chrono 的窄带体积分实现”。

\section{与 PFC 的关系：哪些东西要保留，哪些东西要替换}

\subsection{PFC 中真正本质的结构}
PFC 论文的本质不是“四面体切平面”这一具体实现细节，而是下面这个连续结构：
\begin{enumerate}
    \item 每个物体内部都有一个对象固定的压力场；
    \item 两个物体重叠时，接触面定义为两侧压力相等的地方；
    \item 接触合力和合力矩由该等压面上的 traction 积分给出；
    \item 若不加入耗散，则该模型可以和某种势能解释对应起来。
\end{enumerate}

这四点，才是 PFC 最值得保留的核心。

\subsection{PFC 中可以替换的部分}
PFC 在实现上采用了线性四面体场，并利用“线性场的等压面在 tet--tet 交区中是一张平面”这一性质来做接触面积分。这个实现非常适合其论文当时的网格结构，但它并不是理论本身的一部分。

如果我们改用 SDF 表示几何，那么自然的替代方式是：
\begin{itemize}
    \item 用 SDF 的内部深度来定义压力场；
    \item 用压力差函数的零水平集来定义等压面；
    \item 用 level-set/coarea 公式把面测度改写为体积分测度。
\end{itemize}

因此，从理论上说，\textbf{“SDF 版压力场接触”不是偏离 PFC，而是把 PFC 的连续主干换了一套更适合稀疏体素场的几何表达与数值积分方式。}

\section{第一性出发点：先从一维压缩能量讲清楚}

这一节故意不从三维复杂几何开始，而先从一个外行也能直观看懂的一维压缩问题开始。因为一旦这一维问题讲清楚，后面的三维理论就不再神秘。

\subsection{局部法向柱模型}
把接触区想象成很多根很细很细、彼此独立的“法向小柱子”。对每一根小柱子，两个物体在局部法向上总共关闭了一个距离，记为
\begin{equation}
    \delta \ge 0.
\end{equation}

这个总关闭量由两边分担：
\begin{equation}
    d_A + d_B = \delta,
\end{equation}
其中：
\begin{itemize}
    \item $d_A$ 表示物体 $A$ 在该法向柱上的压缩深度；
    \item $d_B$ 表示物体 $B$ 在该法向柱上的压缩深度。
\end{itemize}

\subsection{每个物体的局部储能与压力律}
对每个物体 $i\in\{A,B\}$，定义单位接触面积上的压缩势能密度
\begin{equation}
    \Psi_i(d_i), \qquad d_i\ge 0.
\end{equation}

对应的法向压力定义为
\begin{equation}
    P_i(d_i)=\Psi_i'(d_i).
\end{equation}

外行理解非常简单：
\begin{itemize}
    \item $\Psi_i(d)$ 表示“压进去 $d$ 这么深，一平方米面积上存了多少弹性能”；
    \item $P_i(d)$ 则表示“再继续压一点点，反推压力会有多大”。
\end{itemize}

\subsection{局部平衡来自最小能量，而不是拍脑袋}
给定总关闭量 $\delta$，局部平衡不应该靠经验猜，而应该由如下极小化问题给出：
\begin{equation}
    \min_{d_A,d_B\ge 0}\; \Psi_A(d_A)+\Psi_B(d_B)
    \quad\text{s.t.}\quad d_A+d_B=\delta.
\end{equation}

这句话的物理意思是：
\begin{quote}
在满足总压缩量固定的前提下，系统会自动选择那种“总能量最省”的压缩分配方式。
\end{quote}

\begin{theorem}[局部平衡条件]
若 $\Psi_A,\Psi_B$ 可微且严格凸，则上述极小化问题有唯一解，且其一阶最优性条件为
\begin{equation}
    P_A(d_A^*) = P_B(d_B^*) = p^*.
\end{equation}
\end{theorem}

\begin{proof}
引入拉格朗日乘子 $\lambda$：
\begin{equation}
    \mathcal L(d_A,d_B,\lambda)=\Psi_A(d_A)+\Psi_B(d_B)+\lambda(d_A+d_B-\delta).
\end{equation}
对 $d_A,d_B$ 求导并令其为零，得到
\begin{equation}
    \Psi_A'(d_A)+\lambda=0,
    \qquad
    \Psi_B'(d_B)+\lambda=0.
\end{equation}
也就是
\begin{equation}
    P_A(d_A)=P_B(d_B)=-\lambda.
\end{equation}
又由于约束 $d_A+d_B=\delta$，故唯一平衡解由“压力相等 + 深度相加等于总关闭量”共同决定。
\end{proof}

这个定理非常重要，因为它说明：\textbf{等压条件不是人为规定，而是局部能量最小化自动推出的。}

\subsection{线性压力律作为第一版模型}
最简单的情形是取
\begin{equation}
    P_i(d)=k_i d,
    \qquad
    \Psi_i(d)=\frac12 k_i d^2.
\end{equation}

这时平衡条件为
\begin{equation}
    k_A d_A = k_B d_B,
    \qquad
    d_A+d_B=\delta.
\end{equation}
解得
\begin{equation}
    d_A^*=\frac{k_B}{k_A+k_B}\,\delta,
    \qquad
    d_B^*=\frac{k_A}{k_A+k_B}\,\delta.
\end{equation}
于是平衡压力为
\begin{equation}
    p^*=k_A d_A^*=k_B d_B^*=\frac{k_Ak_B}{k_A+k_B}\,\delta.
\end{equation}

\begin{corollary}[等效法向刚度]
线性双体接触的等效法向刚度为
\begin{equation}
    k_{\mathrm{eq}}=\frac{k_Ak_B}{k_A+k_B}.
\end{equation}
特别地，当 $k_A=k_B=k$ 时，有
\begin{equation}
    k_{\mathrm{eq}}=\frac{k}{2}.
\end{equation}
\end{corollary}

这条公式对你的 benchmark 校准极其关键。若参考模型采用
\begin{equation}
    F = k_{\mathrm{ref}}\,\delta,
\end{equation}
那么对称双体必须取
\begin{equation}
    k = 2k_{\mathrm{ref}},
\end{equation}
而不是 $k=k_{\mathrm{ref}}$。

\section{从一维局部平衡推广到三维几何}

局部柱模型告诉我们：每个局部法向位置都应该满足“压力相等”。接下来要做的是，把“局部法向位置”这件事写成三维几何中的一个面。

\subsection{每个物体的 SDF}
设两个刚体分别记为 $A$ 与 $B$。在世界坐标下，各自的 signed distance function 记为
\begin{equation}
    \phi_A(x), \qquad \phi_B(x), \qquad x\in\mathbb R^3.
\end{equation}
我们采用常见约定：
\begin{equation}
    \phi_i(x)<0 \text{ 在物体内部},
    \qquad
    \phi_i(x)=0 \text{ 在物体表面},
    \qquad
    \phi_i(x)>0 \text{ 在物体外部}.
\end{equation}

若 $\phi_i$ 在近表面是良好的 signed distance，则有
\begin{equation}
    \norm{\nabla \phi_i(x)}\approx 1,
\end{equation}
且 $-\phi_i(x)$ 可解释为“点 $x$ 深入物体内部的法向距离”。

\subsection{由 SDF 定义内部深度}
定义有效内部深度
\begin{equation}
    d_i(x)=\operatorname{clamp}\bigl(-\phi_i(x),0,H_i\bigr),
\end{equation}
其中 $H_i>0$ 表示该物体可用于构造压力场的有效柔顺层厚度。

之所以使用截断，是因为实现中通常不希望压力无限增大，也不希望窄带之外再继续参与接触计算。

\subsection{对象固定压力场}
对每个物体定义对象固定压力场
\begin{equation}
    p_i(x)=P_i(d_i(x)).
\end{equation}
第一版建议直接取线性形式
\begin{equation}
    p_i(x)=k_i d_i(x),
\end{equation}
因为它最容易校准、最容易解释，也最容易和 benchmark 做一一对应。

\begin{remark}
这一点与 PFC 的核心思想完全一致：物体内部先有一个“对象固定”的压力标量场。不同之处仅在于，PFC 原文用的是基于 tet 的 penetration extent，而这里改成了基于 SDF 的内部深度。
\end{remark}

\section{平衡接触面的严格定义}

\subsection{压力差函数}
定义压力差函数
\begin{equation}
    h(x)=p_A(x)-p_B(x).
\end{equation}

那么平衡接触面自然定义为
\begin{equation}
    \Sigma=\bigl\{x\in\mathbb R^3\mid \phi_A(x)\le 0,\;\phi_B(x)\le 0,\;h(x)=0\bigr\}.
\end{equation}

外行翻译如下：
\begin{itemize}
    \item 点必须在双方的内部重叠区里；
    \item 同时，该点从 $A$ 侧看到的反压和从 $B$ 侧看到的反压要相等；
    \item 满足这两条的点，连起来就是平衡接触面。
\end{itemize}

\subsection{为什么这张面是合理的}
若某点处 $p_A>p_B$，说明站在该点看，$A$ 侧“更抗压”，平衡分界就应该往 $B$ 一侧退；若 $p_A<p_B$，分界面就应向 $A$ 一侧退。只有 $p_A=p_B$ 时，局部法向柱的两边反力平衡。

因此，等压面不是一个视觉上的“中间层”，而是局部法向平衡真正发生的地方。

\subsection{法向由压力差函数给出}
由于 $\Sigma$ 是 $h(x)=0$ 的零水平面，其法向自然为
\begin{equation}
    n(x)=\frac{\nabla h(x)}{\norm{\nabla h(x)}}.
\end{equation}

在线性非饱和区内，若 $0<d_i(x)<H_i$，则
\begin{equation}
    \nabla d_i(x)=-\nabla \phi_i(x),
\end{equation}
从而
\begin{equation}
    \nabla h(x)=-k_A\nabla\phi_A(x)+k_B\nabla\phi_B(x).
\end{equation}

这一公式有一个非常重要的结论：
\begin{quote}
\textbf{接触法向不是简单取某一侧表面法向，也不是把两边法向生硬平均，而是应该取压力差函数的法向。}
\end{quote}

这正是许多数值实现容易悄悄出错的地方。

\section{局部 traction 与总 wrench}

\subsection{无摩擦、无耗散的最小模型}
在最基本的模型中，局部 traction 只包含法向项：
\begin{equation}
    t(x)=p^*(x)\,n(x),
    \qquad x\in\Sigma,
\end{equation}
其中
\begin{equation}
    p^*(x)=p_A(x)=p_B(x)\quad \text{在 }\Sigma\text{ 上成立}.
\end{equation}

于是总合力为
\begin{equation}
    F = \int_{\Sigma} t(x)\,dA,
\end{equation}
相对参考点 $x_O$ 的总合力矩为
\begin{equation}
    \tau_O = \int_{\Sigma}(x-x_O)\times t(x)\,dA.
\end{equation}

\subsection{加入耗散}
若考虑法向耗散，可定义相对速度
\begin{equation}
    v_{AB}(x)=v_A(x)-v_B(x),
\end{equation}
其法向分量为
\begin{equation}
    v_n(x)=n(x)\cdot v_{AB}(x).
\end{equation}

一种稳妥的正定耗散写法是
\begin{equation}
    t_d(x)=\gamma_n\,p^*(x)\,\max\{0,-v_n(x)\}\,n(x),
\end{equation}
其中 $\gamma_n\ge 0$ 为法向耗散系数。

\subsection{加入摩擦}
切向相对速度定义为
\begin{equation}
    v_t(x)=\bigl(I-n(x)n(x)^\top\bigr)v_{AB}(x).
\end{equation}
采用平滑正则化的库仑摩擦，可写为
\begin{equation}
    t_t(x)= -\mu p^*(x)\,\frac{v_t(x)}{\sqrt{\norm{v_t(x)}^2+v_*^2}},
\end{equation}
其中 $\mu$ 为摩擦系数，$v_*>0$ 为正则化速度尺度。

于是总 traction 为
\begin{equation}
    t(x)=p^*(x)n(x)+t_d(x)+t_t(x).
\end{equation}

\section{从面积分到体积分：为什么 band 方案是合法的}

这一步是整套 SDF 路线最关键的理论支点。

\subsection{零水平集测度公式}
设 $\Sigma=\{x\mid h(x)=0\}$，且 $h$ 在 $\Sigma$ 邻域内足够光滑并满足 $\nabla h\neq 0$。对任意足够光滑的测试函数 $\psi(x)$，有
\begin{equation}
    \int_{\Sigma} \psi(x)\,dA
    =
    \int_{\mathbb R^3}
    \psi(x)\,\delta\bigl(h(x)\bigr)\,\norm{\nabla h(x)}\,dV.
\end{equation}

如果还要求点必须落在双方内部重叠区，则再乘以指示函数
\begin{equation}
    \chi_i(x)=
    \begin{cases}
        1,& \phi_i(x)\le 0,\\
        0,& \phi_i(x)>0.
    \end{cases}
\end{equation}
于是
\begin{equation}
    \int_{\Sigma}\psi(x)\,dA
    =
    \int_{\mathbb R^3}
    \psi(x)\,\delta\bigl(h(x)\bigr)\,\norm{\nabla h(x)}\,\chi_A(x)\chi_B(x)\,dV.
\end{equation}

\subsection{应用到接触合力与合力矩}
取
\begin{equation}
    \psi_F(x)=t(x),
    \qquad
    \psi_\tau(x)=(x-x_O)\times t(x),
\end{equation}
即可得到
\begin{equation}
    F = \int_{\mathbb R^3} t(x)\,\delta(h(x))\,\norm{\nabla h(x)}\,\chi_A\chi_B\,dV,
\end{equation}
\begin{equation}
    \tau_O = \int_{\mathbb R^3}(x-x_O)\times t(x)\,\delta(h(x))\,\norm{\nabla h(x)}\,\chi_A\chi_B\,dV.
\end{equation}

这说明：\textbf{厚带体积分不是经验 trick，而是零厚度面积分在 level-set 框架下的严格重写。}

\subsection{数值上为什么需要厚化核}
理想的 Dirac 函数 $\delta$ 无法直接离散实现，因此数值上必须改成一个有界支撑的正则化核 $\delta_\eta$。例如常见的余弦核：
\begin{equation}
    \delta_\eta(s)=
    \begin{cases}
        \dfrac{1}{2\eta}\left(1+\cos\dfrac{\pi s}{\eta}\right), & |s|\le \eta,\\[2mm]
        0, & |s|>\eta.
    \end{cases}
\end{equation}
并要求
\begin{equation}
    \int_{-\infty}^{\infty}\delta_\eta(s)\,ds = 1.
\end{equation}

于是数值积分形式为
\begin{equation}
    F_\eta = \int t(x)\,\delta_\eta(h(x))\,\norm{\nabla h(x)}\,\chi_A\chi_B\,dV,
\end{equation}
\begin{equation}
    \tau_{O,\eta} = \int (x-x_O)\times t(x)\,\delta_\eta(h(x))\,\norm{\nabla h(x)}\,\chi_A\chi_B\,dV.
\end{equation}

\begin{proposition}[窄带极限的一致性]
若 $h\in C^1$，$\nabla h\neq 0$ 于 $\Sigma$ 邻域内，$t$ 连续，且 $\delta_\eta$ 是一族归一化的 Dirac 近似核，则当 $\eta\to 0$ 时，
\begin{equation}
    F_\eta \to F,
    \qquad
    \tau_{O,\eta}\to \tau_O.
\end{equation}
\end{proposition}

\begin{proof}[证明思路]
这本质上是 coarea 公式与近似恒等算子的组合结论。$\delta_\eta$ 把积分限制在 $h=0$ 的一个厚度为 $O(\eta)$ 的邻域内；归一化保证总测度不变；$\eta\to 0$ 时，厚带体积分收缩回真实零水平集面积分。
\end{proof}

\section{Band mechanics 与 Sheet representation 为什么必须分层}

这部分是你前面几份 supplement 中最重要的认识之一。这里把它正式化。

\subsection{两个不同的数学对象}
在连续理论里，本来就有两个不同对象：
\begin{enumerate}
    \item 用于\textbf{求力}的厚化面测度
    \begin{equation}
        \delta_\eta(h)\,\norm{\nabla h};
    \end{equation}
    \item 用于\textbf{表示几何接触片}的零厚度面
    \begin{equation}
        \Sigma=\{x\mid h(x)=0\}.
    \end{equation}
\end{enumerate}

前者是一个“厚带中的面测度权重”；后者是一张真正的零厚度几何面。它们不是一回事。

\subsection{为什么不能直接把厚带当成接触片}
若把 active cells 或 band support 直接解释成接触区域，就会出现：
\begin{itemize}
    \item 面积天然偏厚；
    \item patch 拓扑受厚带层数和体素台阶影响；
    \item bbox 反映的是厚带包围盒，而不是零厚度接触片的支撑域；
    \item 同一份 band 力学结果，换一种 collapse 逻辑就会得到截然不同的几何解释。
\end{itemize}

这说明厚带更适合承担“稳定求力”的任务，而不适合直接承担“正确描述零厚度接触片”的任务。

\subsection{正式分层}
因此应当把接触计算分成两层：

\paragraph{Band mechanics} 负责：
\begin{itemize}
    \item 基于厚带测度稳定地计算总合力与总合力矩；
    \item 保持核正则化、积分连续性与数值鲁棒性。
\end{itemize}

\paragraph{Sheet representation} 负责：
\begin{itemize}
    \item 从 band 样本恢复局部零厚度 sheet 点；
    \item 聚类得到 patch、面积、中心、bbox、pressure center 等几何/拓扑量；
    \item 只负责“怎么表示接触片”，不负责“反过来决定主力学”。
\end{itemize}

\begin{remark}
最简洁的一句话就是：\textbf{Band 负责算得稳，Sheet 负责说得对。}
\end{remark}

\section{一个更自洽的正式理论版本}

下面给出我建议你最终定版时采用的理论结构。它兼容你已有的补充材料，但把逻辑重新收束成一条更短、更闭合的链。

\subsection{层 1：Band mechanics}
对每个世界点 $x$，定义
\begin{equation}
    d_i(x)=\operatorname{clamp}(-\phi_i(x),0,H_i),
    \qquad
    p_i(x)=P_i(d_i(x)).
\end{equation}
令
\begin{equation}
    h(x)=p_A(x)-p_B(x),
    \qquad
    n(x)=\frac{\nabla h(x)}{\norm{\nabla h(x)}}.
\end{equation}

在窄带体积分中，局部平衡压力可近似取为
\begin{equation}
    \bar p(x)=\frac{p_A(x)+p_B(x)}{2}.
\end{equation}
于是 traction 模型写为
\begin{equation}
    t(x)=\bar p(x)n(x)+t_d(x)+t_t(x).
\end{equation}

引入厚带权重
\begin{equation}
    w_\eta(x)=\delta_\eta\bigl(h(x)\bigr)\,\norm{\nabla h(x)}\,\chi_A(x)\chi_B(x).
\end{equation}
则 band 层直接输出
\begin{equation}
    F = \int t(x)w_\eta(x)\,dV,
\end{equation}
\begin{equation}
    \tau_O = \int (x-x_O)\times t(x)w_\eta(x)\,dV.
\end{equation}

\subsection{层 2：Sheet representation}
对每个 band sample $x_i$，定义局部一阶 sheet 投影点
\begin{equation}
    x_i^{\Sigma}=x_i-\frac{h_i}{\norm{\nabla h_i}^2}\,\nabla h_i.
\end{equation}
将带权样本属性
\begin{equation}
    \Delta A_i,\qquad \Delta f_i,\qquad n_i
\end{equation}
一并附着到该局部 sheet seed 上。

随后在三维空间中按几何邻接聚类，而不是在某张全局投影平面上做整数分桶。对每个 cluster $F_\alpha$，定义 measure-preserving collapse：
\begin{equation}
    A_\alpha = \sum_{i\in F_\alpha}\Delta A_i,
\end{equation}
\begin{equation}
    x_\alpha = \frac{\sum_{i\in F_\alpha}\Delta A_i x_i^{\Sigma}}{\sum_{i\in F_\alpha}\Delta A_i},
\end{equation}
\begin{equation}
    n_\alpha = \frac{\sum_{i\in F_\alpha}\Delta A_i n_i}{\left\|\sum_{i\in F_\alpha}\Delta A_i n_i\right\|}.
\end{equation}

这样得到的 $\{x_\alpha,A_\alpha,n_\alpha\}$ 才可用于 area、patch、bbox、support footprint 和可视化输出。

\section{从竖直 column 到一般三维：local-normal accumulator 的正式推导}

前面给出的 band / sheet 分层和 per-column 思路，在平面对平面基准中可以自然拆成“$x$--$z$ 横向 patch 几何 + $y$ 方向法向积分”。这一节把该结构正式推广到一般三维曲面接触：不再依赖全局固定的 $y$ 方向，而是对每个局部接触片在其\textbf{局部法向}上构造一维 band 积分器，并在面内做二维累加。这样得到的 local-normal accumulator 才是 planar-column 版本在一般三维情形下的真正推广。

\subsection{局部坐标系 $(u,v,s)$ 与接触面邻域}

设平衡接触面 $\Sigma$ 在局部可由二维参数 $(u,v)$ 参数化：
\begin{equation}
    r=r(u,v), \qquad (u,v)\in\Omega\subset\mathbb R^2.
\end{equation}
记切向基为
\begin{equation}
    r_u=\frac{\partial r}{\partial u}, \qquad r_v=\frac{\partial r}{\partial v},
\end{equation}
则单位法向可写为
\begin{equation}
    n(u,v)=\frac{r_u\times r_v}{\norm{r_u\times r_v}}.
\end{equation}

在接触面附近取一条局部法向坐标 $s$，定义法向管邻域映射
\begin{equation}
    \Phi(u,v,s)=r(u,v)+s\,n(u,v).
\end{equation}
这里：
\begin{itemize}
    \item $(u,v)$ 描述接触面的面内位置；
    \item $s=0$ 对应零厚度接触面本身；
    \item $s\in[-\varepsilon,\varepsilon]$ 描述围绕接触面的窄带厚度。
\end{itemize}

这一步的本质是把原本直接定义在三维体积中的 band 改写成“面内二维参数 + 法向一维厚度”的局部坐标表示。

\subsection{局部体积元与 Jacobian 修正}

由变量替换公式，
\begin{equation}
    dV = J(u,v,s)\,du\,dv\,ds,
\end{equation}
其中
\begin{equation}
    J(u,v,s)=\left|\det\frac{\partial \Phi}{\partial(u,v,s)}\right|.
\end{equation}

若记曲面的 shape operator 为 $S(u,v)$，则法向管邻域中的 Jacobian 可以写成
\begin{equation}
    J(u,v,s)=\det\!\bigl(I-sS(u,v)\bigr)\,\norm{r_u\times r_v}.
\end{equation}
进一步若以主曲率 $\kappa_1,\kappa_2$ 表示，则
\begin{equation}
    J(u,v,s)=\bigl(1-s\kappa_1(u,v)\bigr)\bigl(1-s\kappa_2(u,v)\bigr)\,\norm{r_u\times r_v}.
\end{equation}

\begin{remark}
当 band 很薄、局部曲率半径远大于 band 厚度时，可以先采用一阶近似
\begin{equation}
    J(u,v,s)\approx \norm{r_u\times r_v}.
\end{equation}
这对应于“每个局部面元沿法向挤出一个近似直棱柱”的工程近似；若曲率效应显著，再把 $\det(I-sS)$ 的修正乘回来。
\end{remark}

\subsection{在局部坐标中重写 band 体积分}

原始的总合力与总合力矩可写为
\begin{equation}
    F=\int_{\mathcal B}\tilde t(x)\,\delta_\eta\bigl(h(x)\bigr)\,\norm{\nabla h(x)}\,dV,
\end{equation}
\begin{equation}
    \tau_O=\int_{\mathcal B}\bigl(\Pi_\Sigma(x)-x_O\bigr)\times \tilde t(x)\,\delta_\eta\bigl(h(x)\bigr)\,\norm{\nabla h(x)}\,dV,
\end{equation}
其中 $\Pi_\Sigma(x)$ 表示把 band 内点投影回零厚度接触面。代入局部坐标系后，得到
\begin{equation}
    F=
    \int_{\Omega}\int_{-\varepsilon}^{\varepsilon}
    \tilde t(u,v,s)\,
    \delta_\eta\bigl(h(u,v,s)\bigr)\,
    \norm{\nabla h(u,v,s)}\,
    J(u,v,s)\,ds\,du\,dv,
\end{equation}
\begin{equation}
    \tau_O=
    \int_{\Omega}\int_{-\varepsilon}^{\varepsilon}
    \bigl(r(u,v)-x_O\bigr)\times
    \tilde t(u,v,s)\,
    \delta_\eta\bigl(h(u,v,s)\bigr)\,
    \norm{\nabla h(u,v,s)}\,
    J(u,v,s)\,ds\,du\,dv.
\end{equation}

这里特意把力矩臂写成 $r(u,v)-x_O$，而不是 $r(u,v)+s\,n(u,v)-x_O$。原因是：band 只是用来近似面测度；真正的接触作用点仍应放回零厚度 sheet 上。这正是前文“Band 负责算得稳，Sheet 负责说得对”的几何化表达。

\subsection{局部 traction 的定义}

在最简洁的无摩擦、无耗散模型中，定义局部平衡压力近似
\begin{equation}
    \bar p(u,v,s)=\frac{p_A(u,v,s)+p_B(u,v,s)}{2},
\end{equation}
并把法向 traction 取为
\begin{equation}
    \tilde t_n(u,v,s)=\bar p(u,v,s)\,n(u,v).
\end{equation}
注意这里使用的是 sheet 上的法向 $n(u,v)$，而不是 band 厚度内每一点各自的法向方向。对薄 band 近似来说，这是一种非常稳定且容易实现的做法。

若加入法向阻尼，可写成
\begin{equation}
    \tilde t_d(u,v,s)=q_d(u,v,s)\,n(u,v),
\end{equation}
例如
\begin{equation}
    q_d(u,v,s)=C_n\,\bar p(u,v,s)\,\max\bigl(0,-v_n(u,v)\bigr),
\end{equation}
其中 $v_n(u,v)$ 是 sheet 上估计的相对法向速度。

若加入正则化切向摩擦，则可先定义每个局部 sheet 元上的法向标量载荷
\begin{equation}
    q_n(u,v)=
    \int_{-\varepsilon}^{\varepsilon}
    \bar p(u,v,s)\,
    \delta_\eta\bigl(h(u,v,s)\bigr)\,
    \norm{\nabla h(u,v,s)}\,
    J(u,v,s)\,ds,
\end{equation}
再定义
\begin{equation}
    \tilde t_t(u,v)=
    -\mu\,q_n(u,v)\,
    \frac{v_t(u,v)}{\sqrt{\norm{v_t(u,v)}^2+v_*^2}},
\end{equation}
其中 $v_t(u,v)$ 是 sheet 上的相对切向速度，$v_*>0$ 是正则化参数。

\subsection{离散化为局部 sheet 元与法向柱}

现在把接触面离散为若干局部 sheet 元，记为 $\alpha=1,\dots,N_\Sigma$。对每个 sheet 元，记录：
\begin{equation}
    \mathcal C_\alpha=
    \left\{
    A_\alpha,\;
    r_\alpha,\;
    n_\alpha,\;
    t_{1,\alpha},\;
    t_{2,\alpha}
    \right\},
\end{equation}
其中：
\begin{itemize}
    \item $A_\alpha$ 是该局部接触片的面元面积；
    \item $r_\alpha$ 是该面元质心；
    \item $n_\alpha$ 是该面元法向；
    \item $t_{1,\alpha},t_{2,\alpha}$ 是局部切向基。
\end{itemize}

然后沿该面元的局部法向离散出若干法向 cell，记为 $m=1,\dots,M_\alpha$，其法向区间为
\begin{equation}
    s\in[s^-_{\alpha,m},s^+_{\alpha,m}].
\end{equation}
定义局部一维函数
\begin{equation}
    h_\alpha(s)=h(r_\alpha+s\,n_\alpha),\qquad
    \bar p_\alpha(s)=\bar p(r_\alpha+s\,n_\alpha),
\end{equation}
以及 Jacobian
\begin{equation}
    J_\alpha(s)=J(u_\alpha,v_\alpha,s).
\end{equation}

\subsection{每个局部法向 cell 的三个核心积分}

对每个局部法向 cell $(\alpha,m)$，定义三个一维积分：
\begin{equation}
    I_{\alpha,m}^{(F)}=
    \int_{s^-_{\alpha,m}}^{s^+_{\alpha,m}}
    \bar p_\alpha(s)\,
    \delta_\eta\bigl(h_\alpha(s)\bigr)\,
    \abs{h_\alpha'(s)}\,
    J_\alpha(s)\,ds,
\end{equation}
\begin{equation}
    I_{\alpha,m}^{(A)}=
    \int_{s^-_{\alpha,m}}^{s^+_{\alpha,m}}
    \delta_\eta\bigl(h_\alpha(s)\bigr)\,
    \abs{h_\alpha'(s)}\,
    J_\alpha(s)\,ds,
\end{equation}
\begin{equation}
    I_{\alpha,m}^{(s)}=
    \int_{s^-_{\alpha,m}}^{s^+_{\alpha,m}}
    s\,
    \delta_\eta\bigl(h_\alpha(s)\bigr)\,
    \abs{h_\alpha'(s)}\,
    J_\alpha(s)\,ds.
\end{equation}

它们分别对应：
\begin{itemize}
    \item 单位面元上该法向 cell 提供的法向力积分；
    \item 单位面元上该法向 cell 对应的面积测度；
    \item 单位面元上的法向一阶矩。
\end{itemize}

再对同一 sheet 元上的所有法向 cell 求和：
\begin{equation}
    I_{\alpha}^{(F)}=\sum_{m=1}^{M_\alpha}I_{\alpha,m}^{(F)},\qquad
    I_{\alpha}^{(A)}=\sum_{m=1}^{M_\alpha}I_{\alpha,m}^{(A)},\qquad
    I_{\alpha}^{(s)}=\sum_{m=1}^{M_\alpha}I_{\alpha,m}^{(s)}.
\end{equation}

\subsection{local-normal 3D accumulator 的离散形式}

对每个局部 sheet 元，定义法向力
\begin{equation}
    F_\alpha^{(n)}=A_\alpha\,I_\alpha^{(F)}\,n_\alpha.
\end{equation}
若加入切向摩擦，则定义
\begin{equation}
    F_\alpha^{(t)}=A_\alpha\,q_\alpha^{(t)},
\end{equation}
其中
\begin{equation}
    q_\alpha^{(t)}=
    -\mu\,I_\alpha^{(F)}\,
    \frac{v_{t,\alpha}}{\sqrt{\norm{v_{t,\alpha}}^2+v_*^2}}.
\end{equation}
于是总局部力为
\begin{equation}
    F_\alpha = F_\alpha^{(n)}+F_\alpha^{(t)}.
\end{equation}

相对参考点 $x_O$ 的局部力矩写为
\begin{equation}
    \tau_\alpha = (r_\alpha-x_O)\times F_\alpha.
\end{equation}
最后，总合力与总合力矩由所有局部 sheet 元累加得到：
\begin{equation}
    F=\sum_{\alpha=1}^{N_\Sigma}F_\alpha,
    \qquad
    \tau_O=\sum_{\alpha=1}^{N_\Sigma}\tau_\alpha.
\end{equation}

\begin{proposition}[local-normal accumulator 的一致性]
若：
\begin{enumerate}
    \item 局部 sheet 元离散收敛到接触面 $\Sigma$；
    \item 每个法向 cell 上的一维积分 $I_{\alpha,m}^{(F)}$、$I_{\alpha,m}^{(A)}$、$I_{\alpha,m}^{(s)}$ 计算一致；
    \item 带宽参数 $\eta$ 与法向离散同时收敛到正确极限；
\end{enumerate}
则上述局部累加器收敛到连续的 band 体积分，从而也收敛到对应的零厚度 traction 面积分。
\end{proposition}

\begin{proof}
这其实就是三个极限的复合：第一，coarea 公式保证零水平面面积分可以由带权体积分表示；第二，局部坐标变换把体积分拆成“面内二维参数积分 + 法向一维积分”；第三，数值离散把二维参数积分离散成 sheet 元求和，把一维法向积分离散成 cell 上求和。只要各层离散都一致，最终总和即收敛到原连续积分。
\end{proof}

\subsection{平面对平面 per-column 是其特例}

前面你已经验证过的 per-column 版本，其实是上述 local-normal 3D accumulator 在一个特殊情形下的退化：
\begin{itemize}
    \item 接触面是平面；
    \item 法向处处相同，$n=(0,1,0)$；
    \item 局部坐标 $(u,v,s)$ 退化为 $(x,z,y)$；
    \item 曲率为零，故 $J\equiv 1$；
    \item 每个局部 sheet 元都对应一个竖直 column 的底面。
\end{itemize}
因此，
\begin{equation}
    (u,v,s)\;\leadsto\;(x,z,y),
\end{equation}
而你先前脚本中的
\[
\texttt{area fraction + local centroid + analytic y-cell integration}
\]
就是 local-normal 3D accumulator 的平面 benchmark 特例，而不是另一套无关的方法。

\subsection{工程上的数据接口建议}

若把这一理论直接翻译成工程接口，那么每个局部接触片至少需要两类数据：

\paragraph{(1) 面内几何数据}
\begin{equation}
    \mathcal C_\alpha=
    \left\{
    A_\alpha,\;
    r_\alpha,\;
    n_\alpha,\;
    t_{1,\alpha},\;
    t_{2,\alpha}
    \right\}.
\end{equation}

\paragraph{(2) 法向积分数据}
\begin{equation}
    \mathcal N_{\alpha,m}=
    \left\{
    s^-_{\alpha,m},\;
    s^+_{\alpha,m},\;
    h_\alpha(s),\;
    \bar p_\alpha(s),\;
    h_\alpha'(s),\;
    J_\alpha(s)
    \right\}.
\end{equation}

然后在代码层面可分成三步：
\begin{enumerate}
    \item 先由几何层构造局部 sheet 元 $\mathcal C_\alpha$；
    \item 再由法向积分器计算各个 $I_{\alpha,m}^{(\cdot)}$；
    \item 最后由 accumulator 负责把所有局部力与力矩加总。
\end{enumerate}

如果以后你要把当前的 planar-column Python 原型升级成一般三维实现，那么最自然的替换关系就是：
\begin{equation}
    y \;\longrightarrow\; s,
\end{equation}
即把“全局竖直方向积分”替换成“每个局部接触片法向上的一维积分”。


\section{正确性应该如何证明}

这里故意不追求“对一切真实弹性现象完全精确”。那样的目标既不现实，也不是这类模型的定位。更合理的正确性分成四层。

\subsection{第一层：局部构成律正确}
如果你把压力律 $P_i(d)$ 视为从局部势能 $\Psi_i(d)$ 导出的导数，那么“等压条件”就是局部能量极小的必要条件。这保证了构成层是自洽的，而不是任意拼装的。

\subsection{第二层：几何定义正确}
若平衡面定义为
\begin{equation}
    \Sigma=\{x\mid \phi_A\le 0,\;\phi_B\le 0,\;p_A=p_B\},
\end{equation}
那么这一定义与 PFC 的 equal-pressure surface 同型。它不会依赖某一个单侧表面法向，也不要求接触区域是单连通的。

\subsection{第三层：面积分到体积分的改写正确}
只要 $h$ 足够光滑且零水平面正规，coarea/level-set 恒等式就保证
\begin{equation}
    \int_{\Sigma}\psi\,dA
    =
    \int \psi\,\delta(h)\,\norm{\nabla h}\,dV.
\end{equation}
因此 band mechanics 有坚实的几何测度基础，而不是启发式近似。

\subsection{第四层：离散实现退化到正确 benchmark}
真正落地时，最关键的是下面三类一致性测试。

\subsubsection{平面一致性}
对平面对平面压缩，接触面积应保持常值，法向合力应严格退化为
\begin{equation}
    F = A\,\frac{k_Ak_B}{k_A+k_B}\,\delta.
\end{equation}

\begin{theorem}[平面一致性]
设两无限平面平行压缩，线性压力律为 $p_i=k_i d_i$，且接触面积为常值 $A$。则 SDF 压力场理论严格退化为
\begin{equation}
    F = A\,\frac{k_Ak_B}{k_A+k_B}\,\delta.
\end{equation}
\end{theorem}

\begin{proof}
平面情形下，局部法向在全域一致，且每一点的总关闭量都等于同一个 $\delta$。于是局部柱模型在整个接触面积上逐点成立，局部平衡压力为
\begin{equation}
    p^*=\frac{k_Ak_B}{k_A+k_B}\,\delta.
\end{equation}
由于 $p^*$ 常值，故总力为
\begin{equation}
    F=\int_A p^*\,dA = A p^* = A\,\frac{k_Ak_B}{k_A+k_B}\,\delta.
\end{equation}
\end{proof}

这个定理告诉我们：\textbf{你的实现必须先过平面 benchmark。若这个 benchmark 都对不上，那么后面一切复杂几何结果都不该先解释成“曲率太复杂”。}

\subsubsection{一阶矩一致性}
对偏心加载，除了总力正确以外，力的作用线还必须正确。这要求至少满足一阶矩一致性：
\begin{equation}
    \tau_O = \int_{\Sigma}(x-x_O)\times t(x)\,dA.
\end{equation}

若数值上 pressure center 已较准但 $F_y$ 仍偏差大，说明问题更像在 band 测度、压力标度或 cell 内积分精度，而不只是几何中心位置错了。

\subsubsection{条带一致性}
对圆柱条带压平面这类 benchmark，应满足：
\begin{itemize}
    \item patch 数应为 1；
    \item 条带宽度应连续增长；
    \item bbox 与 support footprint 不应依赖于 collapse mode 的偶然选择。
\end{itemize}

若同一份 band result 在不同 sheet collapse 逻辑下给出截然不同的 patch 拓扑，那么错误首先是 sheet representation 层，而不是底层物理主线。

\section{为什么你现在的结果可能会和理论对不上}

这里按重要性排序，给出最可能的几类原因。

\subsection{原因一：band 和 sheet 被混用了}
若你直接拿 active cells 或 band region 去解释 area、patch、bbox，那么几何量天然会偏厚。此时就算总力较稳，接触片几何语义也已经漂了。

\subsection{原因二：平面标度没有先校准}
最常见也最容易忽略的一点就是：对称双体线性压力律下，等效刚度不是 $k$，而是 $k/2$。若参考 benchmark 的线性 penalty 刚度是 $k_{\mathrm{ref}}$，则单体 slope 必须先设成 $2k_{\mathrm{ref}}$。

\subsection{原因三：Dirac 厚化核没有正确归一化}
若 $\delta_\eta$ 不满足
\begin{equation}
    \int \delta_\eta(s)\,ds = 1,
\end{equation}
或者 $\eta$ 选得过大、过小，都可能把本应是“面测度”的东西扭曲成“某种厚带体积权重”。这会直接造成法向力和面积系统偏差。

\subsection{原因四：一个 cell 只用中心点采样}
这是你 supplement 3 中已经指出的关键问题。真实平衡面通常是穿过 cell 的，而不是刚好经过 cell 中心。若一个 cell 只用一次中心点评估，就会在：
\begin{itemize}
    \item 边界 cell；
    \item 偏心加载；
    \item 斜穿过格子的薄条带接触
\end{itemize}
中产生系统误差。

\subsection{原因五：SDF 不是好的 signed distance}
若近表面 $\norm{\nabla \phi}\neq 1$，那么 $-\phi$ 就不再是真正的法向距离，而只是某个 level-set 标量。这样会同时污染内部深度、压力场和 $\norm{\nabla h}$ 的测度因子。

\subsection{原因六：你对标的基准可能更接近点接触，而模型本身更擅长 conforming contact}
PFC 原文就明确指出：该类模型在更贴合的接触中更可靠，在更像点接触的情形下会因为未包含 Poisson 效应、剪切效应等而偏软。因此，若 benchmark 本身非常接近 Hertz 型尖锐点接触，那么“并非所有偏差都来自实现 bug”。

\section{建议采用的数值实现闭环}

这一节把上面的理论压缩成一条工程上真正能落地的流程。

\subsection{第 1 步：离线准备}
对每个物体：
\begin{enumerate}
    \item 生成近表面质量足够好的窄带 SDF；
    \item 确保近表面尽量满足 $\norm{\nabla \phi}\approx 1$；
    \item 设定 $H_i$ 和线性 pressure slope $k_i$；
    \item 构建适用于 active leaf / block 的 BVH 或 OBB-BVH。
\end{enumerate}

\subsection{第 2 步：候选对粗检测}
先做刚体级 broadphase，再做 block 级 BVH 遍历，得到所有可能有重叠窄带支持的 block 对。

\subsection{第 3 步：band mechanics 求解}
在每个 block 对的公共积分格内，对每个积分 sample：
\begin{enumerate}
    \item 查询 $\phi_A,\phi_B$ 及其梯度；
    \item 计算 $d_A,d_B,p_A,p_B$；
    \item 计算 $h$、$\nabla h$、$n$；
    \item 计算局部压力近似 $\bar p$；
    \item 计算 $t$ 与厚带权重 $w_\eta$；
    \item 累加
    \begin{equation}
        \Delta F = t\,w_\eta\,\Delta V,
        \qquad
        \Delta \tau_O = (x-x_O)\times t\,w_\eta\,\Delta V.
    \end{equation}
\end{enumerate}

\subsection{第 3.5 步：由全局 sample 组装 local-normal accumulator}
若接触面近似平坦，则可沿全局坐标轴直接做 per-column 累加；而对一般三维接触，更建议把上一步得到的 band sample 先组织为局部 sheet 元，再沿各自法向做一维法向积分。也就是说，工程上应逐步从
\[
\text{全局竖直 column}
\]
过渡到
\[
\text{局部法向 column / prism}.
\]
其最小实现单元应至少包含：
\begin{itemize}
    \item 面内几何量：$A_\alpha,r_\alpha,n_\alpha$；
    \item 法向积分量：$I_\alpha^{(F)},I_\alpha^{(A)},I_\alpha^{(s)}$；
    \item 累加公式：
    \begin{equation}
        F_\alpha = A_\alpha I_\alpha^{(F)} n_\alpha + A_\alpha q_\alpha^{(t)},
        \qquad
        \tau_\alpha = (r_\alpha-x_O)\times F_\alpha.
    \end{equation}
\end{itemize}
这样做的好处是：面内几何误差、法向 band 积分误差与力矩臂误差被彻底分层，可以分别校正。

\subsection{第 4 步：sheet representation 恢复}
对每个 band sample：
\begin{enumerate}
    \item 投影到局部零水平面得到 $x_i^{\Sigma}$；
    \item 携带 $\Delta A_i, \Delta f_i, n_i$ 一起进入 seed 集；
    \item 在世界空间做 fiber / patch 几何邻接聚类；
    \item 做 measure-preserving collapse；
    \item 输出 area、patch count、bbox、sheet center、pressure center 等几何量。
\end{enumerate}

\subsection{第 5 步：校准与验收顺序}
强烈建议严格按下面顺序验收：
\begin{enumerate}
    \item \textbf{平面一致性}：先把平面对平面 benchmark 打通；
    \item \textbf{一阶矩一致性}：再检查偏心压缩的 pressure center 和 $T_z$；
    \item \textbf{条带一致性}：最后检查圆柱条带和平面工况的 patch 与 bbox。
\end{enumerate}

在这三步之前，不建议过早做大规模分辨率 sweep。因为在理论闭环尚未建立时，单纯扫分辨率通常无法告诉你主因究竟是标度问题、测度问题还是表示问题。

\section{你现在最值得优先做的修正}

若目标是尽快让结果与理论重新对齐，优先级建议如下。

\subsection{优先级 1：只保留最小模型}
先关掉摩擦、阻尼和复杂 patch reduction，只保留：
\begin{equation}
    p_i=k_i d_i,
    \qquad
    F = \int \bar p\,n\,\delta_\eta(h)\,\norm{\nabla h}\,dV.
\end{equation}

\subsection{优先级 2：先做平面标定}
用解析关系
\begin{equation}
    k_{\mathrm{eq}}=\frac{k_Ak_B}{k_A+k_B}
\end{equation}
对 benchmark 进行一性标定，不要依靠经验因子慢慢调。

\subsection{优先级 3：把单点采样升级为 cell-wise / sheet-wise quadrature}
至少要让每个活跃 cell 或每个 clipped footprint polygon 拥有多个积分点，而不是“一个 cell 只看一个中心点”或“一个 clipped patch 只看一个质心”。更具体地说，应尽快从
\[
	ext{one cell} 	o 	ext{one point}
\]
升级到
\[
	ext{one clipped polygon} 	o 	ext{triangulation} 	o 	ext{multiple quadrature points} 	o 	ext{multiple local root solves}.
\]
这一项通常会同时改善面积、力和力矩，尤其是在浅嵌入与小接触片工况下。

\subsection{优先级 4：把 band 和 sheet 完全分家}
band 只输出主 wrench；sheet 只负责接触片几何语义。不要再让 sheet 的 patch 构造反过来主导 band 力学。

\subsection{优先级 5：最后再做分辨率扫描}
等前四项完成，再做
\begin{equation}
    \Delta x: 0.05 \to 0.025 \to 0.01 \to 0.005
\end{equation}
这样的分辨率 sweep，才有意义。

\section{总结：最后应当把理论收束成什么样}

本文最终收束出的理论链如下：
\begin{center}
\fbox{SDF 几何 $\to$ 内部深度 $d_i$ $\to$ 压力场 $p_i$ $\to$ 平衡面 $h=0$ $\to$ traction 面积分 $\to$ coarea 窄带体积分 $\to$ band mechanics / sheet representation 分层 $\to$ benchmark 一致性校准}
\end{center}

若进一步面向一般三维接触实现，则其离散层还应继续细化为
\begin{center}
\fbox{sheet 元离散 $\to$ 局部坐标 $(u,v,s)$ $\to$ 法向一维积分 $\to$ local-normal 3D accumulator $\to$ 总 wrench}
\end{center}
这说明你前面在平面对平面基准里验证过的 per-column 方案，并不是另一套独立方法，而只是这一一般三维结构在 $(u,v,s)\mapsto(x,z,y)$ 情形下的退化特例。

如果把这条链每一环都闭合，那么：
\begin{itemize}
    \item 用 SDF 实现 pressure-field contact 不仅可行，而且非常自然；
    \item 它比直接照搬 tet 切面更适合 OpenVDB/NanoVDB 这类稀疏体素场；
    \item 结果对不上时，也能明确知道该先查构成标度、厚带测度、cell 内积分，还是 sheet 表示层。
\end{itemize}

因此，真正需要推翻的不是“SDF 路线”，而是“把 band、sheet、标定和离散测度混成一层”的旧组织方式。只要把这几层分清楚，并先通过平面一致性、一阶矩一致性和条带一致性三类 benchmark，这套理论就是可自洽、可实现、可扩展的。

\appendix
\section{一个建议采用的线性基准版公式清单}

为了便于实现，下面把第一版最小可用模型的公式集中列一次。

\subsection*{几何与深度}
\begin{equation}
    d_i(x)=\operatorname{clamp}\bigl(-\phi_i(x),0,H_i\bigr).
\end{equation}

\subsection*{线性压力场}
\begin{equation}
    p_i(x)=k_i d_i(x).
\end{equation}

\subsection*{平衡函数与法向}
\begin{equation}
    h(x)=p_A(x)-p_B(x),
    \qquad
    n(x)=\frac{\nabla h(x)}{\norm{\nabla h(x)}}.
\end{equation}

\subsection*{局部压力近似}
\begin{equation}
    \bar p(x)=\frac{p_A(x)+p_B(x)}{2}.
\end{equation}

\subsection*{法向 traction}
\begin{equation}
    t_n(x)=\bar p(x)n(x).
\end{equation}

\subsection*{法向耗散}
\begin{equation}
    t_d(x)=\gamma_n\bar p(x)\max\{0,-v_n(x)\}n(x).
\end{equation}

\subsection*{切向摩擦}
\begin{equation}
    t_t(x)= -\mu\bar p(x)\frac{v_t(x)}{\sqrt{\norm{v_t(x)}^2+v_*^2}}.
\end{equation}

\subsection*{厚带权重}
\begin{equation}
    w_\eta(x)=\delta_\eta(h(x))\norm{\nabla h(x)}\chi_A(x)\chi_B(x).
\end{equation}

\subsection*{总力与总力矩}
\begin{equation}
    F = \int t(x)w_\eta(x)\,dV,
    \qquad
    \tau_O = \int (x-x_O)\times t(x)w_\eta(x)\,dV.
\end{equation}

\section{参考文献}
\begin{thebibliography}{99}
\bibitem{elandt2019}
Ryan Elandt, Evan Drumwright, Michael Sherman, Andy Ruina.
\newblock A pressure field model for fast, robust approximation of net contact force and moment between nominally rigid objects.
\newblock arXiv:1904.11433, 2019.

\bibitem{sup1}
补充材料一：基于 OpenVDB/NanoVDB 稀疏窄带 SDF 的体积分压力场接触模型及其在 Chrono 中的实现流程.

\bibitem{sup2}
Supplement 2：基于 OpenVDB/NanoVDB 与 Chrono 的稀疏窄带 SDF 体积分接触算法完整流程说明.

\bibitem{sup3}
Supplement 3：基于 Cell-wise Quadrature 的窄带 SDF 体积分接触修正理论.

\bibitem{sup4}
Supplement 4：双表示接触框架：Band Mechanics 与 Sheet Representation.

\bibitem{sup5}
Supplement 5：Step 6 v2：基于局部 Sheet 点与守测度聚类的接触片表示方法.

\bibitem{sup6}
Supplement 6：当前误差归因、修复路线与可推广算法流程.
\end{thebibliography}

\end{document}
